\chapter{Tiny210 上板调试记录：BL1、SD、USB
与串口}\label{tiny210-ux4e0aux677fux8c03ux8bd5ux8bb0ux5f55bl1sdusb-ux4e0eux4e32ux53e3}

本文记录在 FriendlyARM Tiny210（S5PV210）实板上调通 BL1 → xv6 启动链、
SD 卡 FAT32、EHCI/MT7601U 时遇到的问题、原因和解决办法。每一条都对应
一次实际的串口现象，便于以后遇到同类问题时对照。

相关代码：\passthrough{\lstinline!kernel/arch/arm/mach-s5pv210/bl1/!}（BL1）、
\passthrough{\lstinline!kernel/drivers/mmc/host/!}（SD）、\passthrough{\lstinline!kernel/drivers/usb/host/ehci-hcd.c!}、
\passthrough{\lstinline!ports/tiny210/!}（构建与烧录）。

\section{启动链与各阶段运行在哪里}\label{ux542fux52a8ux94feux4e0eux5404ux9636ux6bb5ux8fd0ux884cux5728ux54eaux91cc}

\begin{lstlisting}
上电 → iROM（BL0，片内 ROM 0xD0000000）
     → 从 SD 第 1 扇区读 16 KiB 到 iRAM 0xD0020000，校验 16 字节头
     → 跳到 0xD0020010：BL1（在 iRAM 中运行）
         时钟 / UART / DDR 初始化 → 读 MBR → FAT32 → xv6_tiny210.img
         → 读到 DDR 0x20007FC0（uImage 数据正好落在 0x20008000）
     → 跳到 0x20008000：xv6 内核（在 DDR 中运行）
         SDHCI → fat32init（FS.IMG 作根文件系统）→ init 挂载 /boot
\end{lstlisting}

\begin{itemize}
\item
  iROM 跳到 BL1 时 \textbf{DDR 还没有初始化}，能用的只有 iROM（只读）和
  iRAM（0xD0020000--0xD0037FFF，96 KiB 片内 SRAM）。所以 BL1 只能在 iRAM
  里运行，最大 16 KiB（含 16 字节头）。
\item
  BL1
  的地址分配（\passthrough{\lstinline!bl1/bl1.ld!}、\passthrough{\lstinline!bl1/start.S!}）：

  {\def\LTcaptype{none} % do not increment counter
  \begin{longtable}[]{@{}p{0.34\linewidth}p{0.60\linewidth}@{}}
  \toprule\noalign{}
  内容 & 地址 \\
  \midrule\noalign{}
  \endhead
  \bottomrule\noalign{}
  \endlastfoot
  代码入口 \passthrough{\lstinline!\_start!} & 0xD0020010（iROM
  固定跳转地址） \\
  代码、常量、数据、BSS & 紧随其后，约 13 KiB \\
  异常向量表（VBAR） & 单独的 \passthrough{\lstinline!.vectors!} 段，32
  字节对齐 \\
  栈 & 0xD0036000 向下 \\
  iROM 自用数据 & 0xD0037xxx（启动通道变量 0xD0037488、拷贝函数指针
  0xD0037F98） \\
  \end{longtable}
  }
\item
  DDR 初始化好以后，大块数据都放在 DDR：内核 0x20007FC0、扇区缓冲
  0x21000000、ATAG 0x20000100。
\end{itemize}

\section{串口：板子正常，终端显示乱码}\label{ux4e32ux53e3ux677fux5b50ux6b63ux5e38ux7ec8ux7aefux663eux793aux4e71ux7801}

\textbf{现象}：\passthrough{\lstinline!minicom!} 里每行只有 7
个左右固定的乱码，如 \passthrough{\lstinline!`J\&?C??!}。

\textbf{原因}：

\begin{enumerate}
\def\labelenumi{\arabic{enumi}.}
\tightlist
\item
  minicom 实际按 9600 接收，而板子按 115200 发送，速率差 12 倍，正好
  对应``一行约 90 个字符只收到 7 个''。
\item
  后台残留了好几个脱离终端的 minicom（甚至有 9600 的），同时抢读
  同一个串口。新开的 minicom 报
  \passthrough{\lstinline!Failed to lock!}，screen 报
  \passthrough{\lstinline!could not find a pty!}。
\end{enumerate}

\textbf{验证方法}：关掉所有 minicom 后，直接用 Python termios 在 115200
读
\passthrough{\lstinline!/dev/cu.usbserial-1110!}，能读到完整文字；改用
9600 读，得到的正是同样的乱码。

\textbf{做法}：

\begin{lstlisting}[language=sh]
screen /dev/cu.usbserial-1110 115200      # 退出：Ctrl-A K y
screen -ls ; screen -r <id>                # 接回已有会话
\end{lstlisting}

\passthrough{\lstinline!make -f ports/tiny210/Makefile uarttest!}
生成一个只打印、不加载内核的
BL1（\passthrough{\lstinline!tiny210-uarttest.bin!}），用来单独检查串口链路。

\section{LED 判断 BL1
走到哪一步}\label{led-ux5224ux65ad-bl1-ux8d70ux5230ux54eaux4e00ux6b65}

Tiny210 核心板的 LED1--4 接在 GPJ2{[}0:3{]}，低电平点亮。BL1
每完成一阶段就 多点亮一个：1 = BL1 在运行，2 = PLL 设好，3 = DDR
自检通过， 4 = 内核已读入、即将跳转。出错停下时，当前阶段的 LED 会闪烁。

注意：复位后 GPJ2 是带下拉的输入，四个灯默认就是全亮的，``全亮''本身
不能说明 BL1 在运行。要看开头是否先闪 3 下。

\section{EHCI 上第一个 bulk OUT
永远发不完}\label{ehci-ux4e0aux7b2cux4e00ux4e2a-bulk-out-ux6c38ux8fdcux53d1ux4e0dux5b8c}

\textbf{现象}：控制传输正常，MT7601U 固件上传的第一个 14 KiB bulk OUT
超时。 qTD 状态是 \passthrough{\lstinline!0xa0!} = Active + \textbf{Data
Buffer Error}。

\textbf{原因}：主机控制器发 OUT 包时，从 DRAM
取数据跟不上（欠载），于是不停 重试。当时 CPU 完全不开
cache，每条指令都要去 DRAM 取，总线被占满。

\textbf{解决}（f9bfbc6）：

\begin{itemize}
\tightlist
\item
  Synopsys EHCI 的 \passthrough{\lstinline!INSNREG01!}（cap+0x94）把
  OUT/IN 包缓冲阈值设成整包 512
  字节（\passthrough{\lstinline!0x00800080!}），整包到齐才开始发；
\item
  打开 I-cache 和分支预测（D-cache 仍关，DMA 不用维护缓存；exec 后会 清
  I-cache）。
\end{itemize}

\textbf{顺带修的隐患}：控制器若支持 64 位地址（HCCPARAMS bit0），会读写
qTD/QH 末尾的 \passthrough{\lstinline!buf\_hi[5]!}。软件字段必须放在这 5
个字之后，和 Linux 一样。 Tiny210 实测
\passthrough{\lstinline!hccparams a010!}，不支持 64
位，但布局已按规范留出位置。

\section{iROM 拷贝函数读第 0
扇区卡死}\label{irom-ux62f7ux8d1dux51fdux6570ux8bfbux7b2c-0-ux6247ux533aux5361ux6b7b}

\textbf{现象}：BL1 打印 \passthrough{\lstinline!SD channel 0!}
后不再有输出。

\textbf{原因}：iROM 的
\passthrough{\lstinline!CopySDMMCtoMem(0xD0037F98)!} 读第 49
扇区（原始内核） 一直正常，但读第 0
扇区（MBR）会卡死。最初还把目标放在了 iRAM 的静态 数组里，改到 DDR
后仍卡，说明问题在第 0 扇区本身。

\textbf{解决}（49da8ac）：不再依赖 iROM 读 FAT32。把内核 SD
驱动的协议部分拆成
\passthrough{\lstinline!kernel/drivers/mmc/host/sdcore.c!}（无内核依赖、无除法），BL1
和内核共用：

\begin{itemize}
\tightlist
\item
  时间基准用 PWM 定时器 3 自由计数（PCLK\_PSYS/16）；
\item
  引脚 GPG0、HSMMC0 时钟门控、SCLK\_MMC 取 24 MHz 晶振，和内核一致；
\item
  卡识别 375 kHz，数据 24 MHz、4 位、单块 PIO；
\item
  自己的初始化失败时，退回 iROM 函数（最后一个参数
  \passthrough{\lstinline!init=1!} 让它重新 初始化卡），读第 49
  扇区的后备内核。
\end{itemize}

\section{初始化``无声失败''：R1b 的 busy
结束判断}\label{ux521dux59cbux5316ux65e0ux58f0ux5931ux8d25r1b-ux7684-busy-ux7ed3ux675fux5224ux65ad}

\textbf{现象}：\passthrough{\lstinline!own SD init failed!}，但没有任何具体原因。

\textbf{原因}：CMD2/3/7 失败时原来不打日志。CMD7（SELECT\_CARD）是带
busy 的 R1b 响应，原来只等 Transfer Complete
中断，这个控制器在没有数据阶段时 不一定会产生它。

\textbf{解决}（0cec709）：每一步都打日志；R1b 的 busy 期以 Transfer
Complete \textbf{或} PRNSTS bit20（DAT0
电平）变高为结束。之后初始化成功：
\passthrough{\lstinline!SD 2.0 card, ocr 0xc0ff8000!}、\passthrough{\lstinline!SDHC/SDXC card, 4-bit 24 MHz, rca 0x21e1!}。

\section{没有异常向量表时，出错只会``卡死''}\label{ux6ca1ux6709ux5f02ux5e38ux5411ux91cfux8868ux65f6ux51faux9519ux53eaux4f1aux5361ux6b7b}

正常执行不需要向量表，CPU 只是顺序取指令；只有发生异常（data abort、
未定义指令、中断）时才会跳到``向量基址 + 偏移''。BL1 原来没有设置 VBAR，
异常会进入 iROM 自己的处理程序，表现为无输出卡死或重新启动，看不出原因。

\textbf{解决}（eb41bd3）：BL1 启动时设置 VBAR
指向自己的向量表，异常时打印：

\begin{lstlisting}
BL1: CPU exception: data abort lr=... dfar=... dfsr=... ifsr=...
\end{lstlisting}

内核另有自己的向量表（\passthrough{\lstinline!kernel/arch/arm/vectors.S!}），因为它依赖定时器
中断和 SVC 系统调用。

\section{\texorpdfstring{一行 \texttt{.balign\ 32} 让所有地址错开 16
字节}{一行 .balign 32 让所有地址错开 16 字节}}\label{ux4e00ux884c-balign-32-ux8ba9ux6240ux6709ux5730ux5740ux9519ux5f00-16-ux5b57ux8282}

\textbf{现象}：加上向量表后，从第一行起字符串就成了乱码，并走进各种错误分支，
最后触发异常、复位。

\textbf{原因}：VBAR 要求 32 字节对齐，向量表写在
\passthrough{\lstinline!.text.start!} 里并用了
\passthrough{\lstinline!.balign 32!}。链接器因此把整个
\passthrough{\lstinline!.text!} 对齐到 32 字节，起始地址变成
0xD0020020。但 iROM 总是把镜像放到 0xD0020010 再跳过去，于是所有绝对
地址（字符串、常量池、跳转表）都偏了 16 字节。

\textbf{解决}（fcdbbce）：向量表放进单独的
\passthrough{\lstinline!.bl1vectors!} 段，排在
\passthrough{\lstinline!.text!} 之后； \passthrough{\lstinline!bl1.ld!}
增加断言
\passthrough{\lstinline!ADDR(.text) == 0xD0020010!}，入口地址不对时编译直接报错。

检查办法：

\begin{lstlisting}[language=sh]
arm-none-eabi-objdump -h build/tiny210/bl1/bl1.elf   # .text 必须从 d0020010 开始
xxd -s 16 -l 16 build/tiny210/tiny210-bl1.bin        # 第一条指令 cpsid if = c0 00 0c f1
\end{lstlisting}

\section{解析 MBR
时的对齐错误}\label{ux89e3ux6790-mbr-ux65f6ux7684ux5bf9ux9f50ux9519ux8bef}

\textbf{现象}：

\begin{lstlisting}
BL1: sd: CMD17, lba 0 ... transfer complete
BL1: CPU exception: data abort lr=0xd00217fc dfar=0x210001c6 dfsr=0x00000001
\end{lstlisting}

\textbf{分析}：\passthrough{\lstinline!dfsr=0x1!} 是 alignment
fault；\passthrough{\lstinline!dfar=0x210001c6!} = 扇区缓冲 0x21000000 +
0x1C6，正是 MBR 第一个分区表项的起始 LBA（446+8）， 不是 4
字节对齐的地址。代码是逐字节读取的，但 clang 针对 Cortex-A8 把
它合并成了一条非对齐的 \passthrough{\lstinline!ldr!}；而 iROM 留下的
SCTLR.A=1（开着对齐检查）。

\textbf{解决}（6653ddc）：BL1 用
\passthrough{\lstinline!-mno-unaligned-access!} 编译，并在启动时清
SCTLR.A。内核的 \passthrough{\lstinline!entry.S!} 本来就清了 A 位。

\section{SD 卡布局与文件}\label{sd-ux5361ux5e03ux5c40ux4e0eux6587ux4ef6}

{\def\LTcaptype{none} % do not increment counter
\begin{longtable}[]{@{}p{0.34\linewidth}p{0.60\linewidth}@{}}
\toprule\noalign{}
位置 & 内容 \\
\midrule\noalign{}
\endhead
\bottomrule\noalign{}
\endlastfoot
第 0 扇区 & MBR（不动） \\
第 1--32 扇区 & BL1（16 KiB） \\
第 49 扇区起 & 后备内核（原始 uImage，约 8.6k 扇区） \\
第 31248 扇区起 & FAT32（\passthrough{\lstinline!/Volumes/bootfs!}） \\
FAT32 根目录 &
\passthrough{\lstinline!xv6\_tiny210.img!}（内核）、\passthrough{\lstinline!FS.IMG!}（根文件系统，可写）、\passthrough{\lstinline!mt7601u.bin!} \\
\end{longtable}
}

\begin{itemize}
\tightlist
\item
  \passthrough{\lstinline!xv6\_tiny210.img!} 是 VFAT 长文件名，BL1
  不区分大小写匹配。
\item
  内核镜像里还内嵌了一份 4 MiB 的 fs.img 作为后备根文件系统， 卡上没有
  \passthrough{\lstinline!FS.IMG!} 时使用（只在内存中，修改不保存）。
\item
  \passthrough{\lstinline!make tiny210!} 会把内核和固件拷到
  \passthrough{\lstinline!/Volumes/bootfs!}；\passthrough{\lstinline!FS.IMG!}
  只在卡上 没有时才写（\passthrough{\lstinline!BOOTFS\_FSIMG=1!}
  强制覆盖）。BL1 变化时需要手动：
\end{itemize}

\begin{lstlisting}[language=sh]
sudo dd if=build/tiny210/tiny210-bl1.bin of=/dev/rdisk4 bs=512 seek=1
\end{lstlisting}

\section{固件与文件的加载时机}\label{ux56faux4ef6ux4e0eux6587ux4ef6ux7684ux52a0ux8f7dux65f6ux673a}

\begin{itemize}
\tightlist
\item
  根文件系统挂载之前就要用的东西只能随内核一起带进来：\passthrough{\lstinline!mt7601u.bin!}
  的内置副本用 \passthrough{\lstinline!.incbin!}
  链接进内核（\passthrough{\lstinline!request\_firmware()!} 的后备）。
\item
  根文件系统之后：\passthrough{\lstinline!fat32init()!} 在 initcall
  之后运行，FAT32 卷在内核里 直接可读；init 把它挂到
  \passthrough{\lstinline!/boot!}。\passthrough{\lstinline!request\_firmware()!}
  先读 FAT32 根目录， 失败再用内置副本。
\item
  MT7601U 在 \passthrough{\lstinline!ehci-enum!}
  内核线程里枚举，此时调度器、中断和定时器都已 运行，FAT32 也已就绪。
\end{itemize}

\section{内核启动前的长时间停顿}\label{ux5185ux6838ux542fux52a8ux524dux7684ux957fux65f6ux95f4ux505cux987f}

\textbf{现象}：BL1 打印
\passthrough{\lstinline!starting xv6 at 0x20008000!}
之后，要停好几秒内核才有输出。

\textbf{原因}：初始化页分配器时，xv6 把 128 MiB
空闲内存每一页都填``垃圾值'' （\passthrough{\lstinline!kfree!} 里的
\passthrough{\lstinline!memset(va, 1, PGSIZE)!}，用来发现悬空指针）。内核的
\passthrough{\lstinline!memset!} 是逐字节循环，D-cache
又是关的，相当于对 DRAM 做了约 1.3 亿次 未缓存的单字节写。

\textbf{解决}（e1263cd）：\passthrough{\lstinline!memset!}/\passthrough{\lstinline!memmove!}
改成按对齐的 32 位字操作； \passthrough{\lstinline!freerange()!}
在启动时把从未用过的页交给分配器时跳过填充，之后正常
\passthrough{\lstinline!kfree!} 仍然填充。

\section{Wi-Fi 扫描永远停在信道
1：延时任务从来不执行}\label{wi-fi-ux626bux63cfux6c38ux8fdcux505cux5728ux4fe1ux9053-1ux5ef6ux65f6ux4efbux52a1ux4eceux6765ux4e0dux6267ux884c}

\textbf{现象}：扫描列出的 AP 全都是
\passthrough{\lstinline!channel=1!}；\passthrough{\lstinline!wifi!} 打印
\passthrough{\lstinline!joining TP-Link\_B114!} 后 30 秒超时，始终没有
\passthrough{\lstinline!target ... found!}。

\textbf{原因}：workqueue
的延时任务（\passthrough{\lstinline!queue\_delayed\_work!}）靠
\passthrough{\lstinline!workqueue\_timer\_tick()!} 推进
\passthrough{\lstinline!work\_jiffies!}。树莓派的定时器中断会调用它，
ARM32 移植时漏掉了，\passthrough{\lstinline!work\_jiffies!} 一直是
0。于是 MT7601U 的 \passthrough{\lstinline!state\_work!}
（定时换信道、认证/关联超时重试）和
\passthrough{\lstinline!wifi\_connect!} 安排的``立即开始连接''
都永远不会触发。

为什么系统其他部分完全正常：进程调度、\passthrough{\lstinline!sleep()!}、\passthrough{\lstinline!uptime!}
用的是另一套 计时（\passthrough{\lstinline!ticks!} 和
\passthrough{\lstinline!sched\_tick()!}，在
\passthrough{\lstinline!arm32\_timer\_tick()!} 里推进）；立即执行的
\passthrough{\lstinline!queue\_work()!}（例如 mt7601u 的收包
NAPI）也不依赖 \passthrough{\lstinline!work\_jiffies!}。只有延时
任务受影响，所以问题很隐蔽。

\textbf{解决}（0bd029e）：\passthrough{\lstinline!arm32\_timer\_tick()!}
每个节拍（10 ms，与树莓派相同） 调用一次
\passthrough{\lstinline!workqueue\_timer\_tick()!}。之后扫描覆盖 1--11
信道，在信道 9 找到 TP-Link\_B114，认证、关联成功，有一次 WPA2
四次握手完整完成。

\section{超时后清理 QH 把数据 toggle
复位了}\label{ux8d85ux65f6ux540eux6e05ux7406-qh-ux628aux6570ux636e-toggle-ux590dux4f4dux4e86}

\textbf{现象}：某个 bulk OUT 超时后重新连接，AP 反复发 M1，驱动反复回 M2
（replay 计数递增），AP 始终收不到。

\textbf{原因}：EHCI 驱动在超时后用 \passthrough{\lstinline!qh\_flush()!}
清理 QH，把 overlay 的 token 整个 清零，数据 toggle 也就回到了
DATA0。但设备端的 toggle 没变，于是之后每个 包都被设备当成重发包，ACK
之后又悄悄丢弃。主机这边看来全部``发送成功''。

\textbf{解决}（f649f28）：清理时只清除 Active/Halted/Ping，保留 DT 位；
bulk OUT 的超时从 1 秒放宽到 3 秒，给 PING/NAK 中的设备更多时间。

\section{MT7601U 发送端点一直
NAK（排查中）}\label{mt7601u-ux53d1ux9001ux7aefux70b9ux4e00ux76f4-nakux6392ux67e5ux4e2d}

\textbf{现象}：认证、关联、EAPOL 等前若干帧都能发出去，之后 EP8 的 bulk
OUT 一直不完成。qTD 状态 \passthrough{\lstinline!0x80!}（只有
Active），QH token 带 Ping 位，说明设备对 PING 一直回
NAK，即设备侧发送缓冲没有腾出来。

\textbf{尝试 1：换端点}（da79cff，已撤回）。MT7601U 有 6 个 bulk OUT
（0x08、0x04、0x05、0x06、0x07、0x09）。参照 Linux mt7601u，0x08 是 MCU
带内命令 端点，管理队列（QSEL\_MGMT）对应第 6 个端点 0x09，于是把 802.11
帧改到 EP9。 结果 EP9 从第一个认证帧就 NAK，比 EP8
更糟，说明在本驱动当前的初始化 状态下这个推测不成立。USB core
现在会记录全部 bulk OUT 端点
（\passthrough{\lstinline!usb\_device.bulk\_out\_eps[]!}），方便后续继续试验。

\textbf{尝试 2：读空 TX 状态 FIFO}（5ac8c59，待上板验证）。驱动从未读取
\passthrough{\lstinline!TX\_STAT\_FIFO\_EXT!}（0x1798）/\passthrough{\lstinline!TX\_STAT\_FIFO!}（0x1718），而
Linux 每次发送后都会读。 推测 FIFO 满后 MAC 不再接收新帧，于是端点一直
NAK，与``前几帧正常、之后 永远 NAK''的现象吻合。现在每发一帧就读空
FIFO，并打印前 8 条状态 （\passthrough{\lstinline!acked!} /
\passthrough{\lstinline!not acked!}）用于确认。

\section{经验小结}\label{ux7ecfux9a8cux5c0fux7ed3}

\begin{enumerate}
\def\labelenumi{\arabic{enumi}.}
\tightlist
\item
  先确认观察手段：串口终端设置、后台进程、LED 默认状态，都可能造成误判。
\item
  早期代码一定要有异常向量和逐步日志，否则所有问题看起来都是``卡死''。
\item
  位置相关的代码（BL1）要在链接脚本里断言入口地址。
\item
  iROM / Boot ROM 提供的函数只在它们被验证过的用法下可靠（这里：读 第 49
  扇区到 DDR），超出这个范围就自己实现。
\item
  iROM 留下的 CPU 状态（SCTLR.A、V 位、cache）要显式设成自己需要的值。
\item
  主机测试只能验证逻辑（FAT32 解析、CRC、ATAG），验证不了内存类型、
  对齐、时序这些硬件相关的问题。
\item
  从另一个平台移植代码时，要核对它依赖的``周期性调用''（定时器节拍里的
  \passthrough{\lstinline!workqueue\_timer\_tick()!}
  这类）是否都接上了，缺了往往不会崩溃，只会 让某些功能静静地失效。
\item
  出错恢复路径（超时、flush）本身也要正确：错误恢复时丢掉的协议状态
  （USB 数据 toggle）会让之后所有``成功''的传输变成无声丢包。
\item
  不熟悉的硬件行为（端点用途、状态 FIFO）先对照 Linux 驱动，但改动要能
  上板验证并且可以撤回，像第 15 节那样记录每一次尝试的结果。
\end{enumerate}
